\documentclass[12pt]{article}

% Packages
\usepackage{amsmath}
\usepackage{geometry}
\usepackage{setspace}
\usepackage{mdframed}
\usepackage{framed}
\usepackage{xcolor}
% Page setup
\geometry{margin=1in}
\onehalfspacing

\title{Waste Avoided Estimation Report}


\begin{document}

\maketitle

\section{Waste Avoided Estimation and Assumptions}


\subsection{Estimation Method}
The system tracks food waste avoided by estimating the weight of food redistributed through bundles. Waste avoided is calculated using the following formula:

\begin{equation}
\text{Waste Avoided (kg)} = \text{Number of Collected Bundles} \times \text{Average Weight per Bundle (kg)}
\end{equation}


\subsection{Key Assumptions}
\begin{itemize}

    \item Each bundle contains approximately X kg of food on average.
    \item This estimate is based on typical surplus food bundle sizes from similar food redistribution platforms and internal testing.
    \item The weight is assumed to be constant across bundles for the purposes of analytics.
\end{itemize}




\subsection{These assumptions are}

\begin{itemize}
    \item Hard-coded in the analytics pipeline for consistency across all calculations.
    \item Waste Avoided data  is reflected in the Seller Analytics dashboard.
\end{itemize}




\subsection{Example}

\begin{framed}
\begin{itemize}
    \item Waste avoided is estimated.
    \item Assumption: each collected bundle contains $\sim$1.5 kg of food on average.
    \item Calculation: \( \text{wasteAvoidedKg} = \text{bundlesCollected} \cdot 1.5 \)
\end{itemize}
\end{framed}






\subsection{Summary}
Waste avoided is estimated using the number of collected bundles multiplied by an assumed average bundle weight (e.g., 1.5 kg per bundle)
\end{document}